\documentclass[11pt]{article}

\usepackage[utf8]{inputenc}
\usepackage[T1]{fontenc}
\usepackage{amsmath,amssymb,amsthm,mathtools}
\usepackage{geometry}
\usepackage{hyperref}
\usepackage{enumitem}
\usepackage{booktabs}

\geometry{margin=1in}
\setlength{\emergencystretch}{3em}

\newtheorem{theorem}{Theorem}
\newtheorem{proposition}{Proposition}
\newtheorem{definition}{Definition}
\newtheorem{remark}{Remark}

\newcommand{\C}{\mathbb{C}}
\newcommand{\Q}{\mathbb{Q}}
\newcommand{\Z}{\mathbb{Z}}
\newcommand{\On}{\mathcal{O}}
\newcommand{\Weyl}{A_n}
\newcommand{\dR}{\mathrm{dR}}
\newcommand{\STr}{\mathrm{STr}}
\newcommand{\Conf}{\mathrm{Conf}}

\title{A Finite Rosetta Layer for q-Super-Cuntz Regularization,\\
Configuration Complements, and D-Module de Rham Computation}

\author{Formalization draft generated from the \texttt{/home/goutev/auto} Lean/SymPy repository}
\date{June 18, 2026}

\begin{document}
\maketitle

\begin{abstract}
We report a small formalization layer connecting four finite algebraic
regularization mechanisms: the endogenous Krein signature, the supertrace
pairing, hyperbolic spectral squashing, and Cayley unitary boundary
coordinates.  The root object is a two-sector q-super-Cuntz socket
\(\On_q(1|1)\), represented at finite level by a \(2\times2\) matrix model in
which the superparity operator equals the Krein signature
\(\eta=N_+-N_-\).  The work is formalized in Lean 4 as theorem-honest finite
sockets and checked in SymPy by executable witnesses.  We further add a
three-point Orlik--Solomon skeleton for non-isotropic configuration
complements and a D-module de Rham computation socket based on the
Oaku--Takayama algorithm for complements of affine hypersurfaces.  The paper
does not claim a completed analytic q-to-1 colimit, full root-of-unity
representation theory, or a physical spacetime reconstruction theorem.  Those
claims are explicitly represented as sockets.
\end{abstract}

\section{Introduction}

The repository formalizes a finite algebraic spine for a larger program in
which metric signature, regularization, and complement cohomology are treated
as outputs of algebraic structure rather than as external analytic choices.
The central finite observation is the identity
\[
  \eta = N_+ - N_- = (-1)^F
\]
in a two-sector q-super-Cuntz model.  From this identity follow the finite
supertrace
\[
  \STr(X)=\operatorname{Tr}(\eta X)=X_{00}-X_{11},
\]
the involutivity \(\eta^2=1\), and a bounded/unitary coordinate pipeline
through hyperbolic squashing and the Cayley transform.

This document describes exactly what has been implemented and verified:
finite matrix identities, formal Orlik--Solomon relations, a finite
root-of-unity truncation socket, and a Weyl-algebra witness for the
Oaku--Takayama D-module algorithm.  The goal is not to oversell the analytic
part of the program, but to isolate small objects that can be compiled,
tested, and extended.

\section{The q-Super-Cuntz Regularization Rosetta Layer}

Let
\[
  N_+ =
  \begin{pmatrix}
  1&0\\0&0
  \end{pmatrix},
  \qquad
  N_- =
  \begin{pmatrix}
  0&0\\0&1
  \end{pmatrix},
  \qquad
  \eta=N_+-N_-.
\]
The finite q-super-Cuntz socket uses the two sectors as bosonic and fermionic
projections:
\[
  P_B=N_+,\qquad P_F=N_-,\qquad (-1)^F=P_B-P_F.
\]

\begin{proposition}[Finite superparity signature]
In the \(2\times2\) socket,
\[
  (-1)^F=\eta,\qquad \eta^2=1,\qquad \eta^\dagger=\eta.
\]
\end{proposition}

\begin{proof}
This is a direct matrix calculation.  In the repository the corresponding
Lean module proves the three identities as named theorems:
\[
\texttt{superParity\_eq\_eta},\qquad
\texttt{superParity\_sq},\qquad
\texttt{superParity\_selfadjoint}.
\]
\end{proof}

\begin{proposition}[Finite supertrace]
For any \(2\times2\) matrix \(X\),
\[
  \STr(X)=\operatorname{Tr}((-1)^F X)=X_{00}-X_{11}.
\]
\end{proposition}

\begin{proof}
This is proved in Lean as \path{supertrace_apply}; the SymPy witness
\path{q_super_regularization_rosetta.py} checks the same identity.
\end{proof}

\subsection{The four regularization roles}

The Lean module \texttt{QSuperRegularizationRosetta.lean} packages the finite
facts into four named projections.

\begin{center}
\begin{tabular}{lll}
\toprule
Role & q-super origin & finite contract\\
\midrule
Minkowski signature & \((-1)^F=P_B-P_F\) & \(\eta^2=1\), \(\eta^\dagger=\eta\)\\
Lorentz bilinear & \(\STr(X)=\operatorname{Tr}(\eta X)\) & \(X_{00}-X_{11}\)\\
Spectral squashing & q-integer shadow & \(-1<\tanh(1-\lambda^{-1})<1\)\\
CUE/q-circle & Cayley coordinate & \(U U^\dagger=1\)\\
\bottomrule
\end{tabular}
\end{center}

The analytic q-to-1 colimit, root-of-unity representation theory, and TKK
spacetime recovery are represented as explicit propositions, not as proved
theorems.

\section{Finite Root-of-Unity Truncation}

The module \texttt{QRootOfUnityTruncation.lean} isolates the finite content of
the phrase ``root-of-unity truncation.''  A stage is a natural number \(k>0\),
with state space \(\operatorname{Fin}(k)\).

\begin{definition}
At stage \(k\), a natural level \(n\) is admissible if \(n<k\).
\end{definition}

\begin{proposition}[Hard boundary cutoff]
The first boundary level \(k\) is not admissible:
\[
  \neg(k<k).
\]
The finite state space has exactly \(k\) states.
\end{proposition}

The same module defines a parity sign on the finite ladder and proves
\[
  \operatorname{paritySign}(n)^2=1.
\]
This captures the finite combinatorial part of a root-of-unity cutoff without
claiming the full analytic representation theorem for \(q=\exp(i\pi/k)\).

\section{A Three-Point Orlik--Solomon Complement Skeleton}

For three points in \(\C^D\), consider the non-isotropic complement where
\[
  q(x_i-x_j)\neq 0
\]
for all pairs \(i<j\).  The repository formalizes the finite
Orlik--Solomon skeleton with generators
\[
  A_{12},\quad A_{23},\quad A_{13}
\]
and the Arnold--Orlik--Solomon relation
\[
  A_{12}A_{23}-A_{12}A_{13}+A_{23}A_{13}=0.
\]

\begin{proposition}
After imposing the Arnold--Orlik--Solomon relation, the formal degree-two
space has rank \(2\), with normal basis
\[
  A_{12}A_{23},\qquad A_{12}A_{13}.
\]
\end{proposition}

The Lean theorem \texttt{arnold\_orlik\_solomon\_relation} proves that the
relation reduces to zero in this normal form.  The SymPy witness
\texttt{non\_iso\_conf3\_orlik\_solomon.py} checks the concrete \(D=4\)
quadratic form
\[
  q(z)=z_0^2+z_1^2+z_2^2+z_3^2
\]
and verifies symbolically that the forms \(\mathrm{d}\log q(x_i-x_j)\) are
closed.

\begin{remark}
Analytically, \(\mathrm{d}\log q\) is a one-form on the hypersurface
complement.  The repository keeps the requested \(D-1\) grading as a formal
socket rather than silently identifying it with the analytic de Rham degree.
\end{remark}

\subsection{Cooperad map}

For the collapse \(\{1,2\}|3\), the formal cooperad map is
\[
  \Delta_{12}(A_{12})=1\otimes A_{12},\qquad
  \Delta_{12}(A_{13})=A_{13}\otimes 1,\qquad
  \Delta_{12}(A_{23})=A_{13}\otimes 1.
\]
This is implemented by the Lean function \texttt{delta12}.

\section{D-Module de Rham Computation Socket}

The paper \cite{OakuTakayama1998} gives an algorithm for computing
\[
  H^k(U,\C_U),\qquad U=\C^n\setminus V(f),
\]
for nonzero \(f\in\Q[x_1,\ldots,x_n]\).  The implemented Lean socket
\texttt{OakuTakayamaDModuleDeRham.lean} records the algorithmic pipeline:
\[
  \Q[x,1/f]\simeq A_n/I,
\]
then a formal Fourier transform
\[
  x_i\mapsto -\partial_i,\qquad \partial_i\mapsto x_i,
\]
then a free resolution of \(A_n/J\), and finally the restriction/integration
complex
\[
  A_n/(x_1A_n+\cdots+x_nA_n)\otimes_{A_n} A_n^{p_{-k}}.
\]

\subsection{The introductory Weyl algebra atom}

In the one-variable Weyl algebra \(A_1=\Q\langle x,\partial\rangle\) with
\[
  \partial x=x\partial+1,
\]
the paper considers
\[
  p=(x-u)(x-v)\partial-a(x-v)-b(x-u).
\]
The SymPy witness implements normal ordering and verifies that the formal
Fourier transform is
\[
  \widehat p
  =
  x\partial^2
  +\bigl((u+v)x+2+a+b\bigr)\partial
  +uvx+u+v+av+bu.
\]
The associated introductory b-polynomial atom is
\[
  b(s)=s(s-a-b),
\]
and Lean proves
\[
  b(0)=0,\qquad b(a+b)=0.
\]

\section{Implementation}

\begin{center}
\begin{tabular}{ll}
\toprule
File & Role\\
\midrule
\texttt{QSuperCuntzRegularization.lean} & finite q-super-Cuntz identities\\
\texttt{QSuperRegularizationRosetta.lean} & four regularization projections\\
\texttt{QRootOfUnityTruncation.lean} & finite root-of-unity cutoff socket\\
\texttt{NonIsoConf3OrlikSolomon.lean} & three-point OS/cooperad skeleton\\
\texttt{OakuTakayamaDModuleDeRham.lean} & D-module de Rham algorithm socket\\
\texttt{q\_super\_regularization\_rosetta.py} & finite matrix witness\\
\texttt{non\_iso\_conf3\_orlik\_solomon.py} & D=4 complement witness\\
\texttt{oaku\_takayama\_dmodule\_de\_rham.py} & Weyl \(A_1\) Fourier witness\\
\bottomrule
\end{tabular}
\end{center}

The current verification commands are:
\[
\begin{aligned}
&\texttt{lake build OakuTakayamaDModuleDeRham NonIsoConf3OrlikSolomon},\\
&\texttt{python3 proofs/oaku\_takayama\_dmodule\_de\_rham.py},\\
&\texttt{python3 proofs/non\_iso\_conf3\_orlik\_solomon.py}.
\end{aligned}
\]

\section{Limitations and Next Steps}

The present work is intentionally finite and modular.  It does not yet prove:
\begin{itemize}[leftmargin=2em]
\item the analytic q-to-1 colimit;
\item full root-of-unity representation theory for \(\On_q(1|1)\);
\item an analytic comparison theorem for the non-isotropic configuration
  complement;
\item a Lean implementation of Weyl algebra Gröbner bases;
\item a full TKK reconstruction theorem for macroscopic spacetime.
\end{itemize}

The next mathematically conservative extension is to replace the D-module
algorithm sockets by a computable Weyl algebra normal-form layer, then connect
that layer to small explicit complement examples.

\begin{thebibliography}{9}

\bibitem{OakuTakayama1998}
Toshinori Oaku and Nobuki Takayama.
\newblock An algorithm for de Rham cohomology groups of the complement of an affine variety via D-module computation.
\newblock \emph{arXiv:math/9801114}, 1998.

\bibitem{OrlikTerao}
Peter Orlik and Hiroaki Terao.
\newblock \emph{Arrangements of Hyperplanes}.
\newblock Springer, 1992.

\bibitem{Bjork}
Jan-Erik Bjork.
\newblock \emph{Rings of Differential Operators}.
\newblock North-Holland, 1979.

\bibitem{Deligne}
Pierre Deligne.
\newblock Equations differentielles a points singuliers reguliers.
\newblock \emph{Lecture Notes in Mathematics}, Vol. 163, Springer, 1970.

\end{thebibliography}

\end{document}
